\documentclass[12pt,a4paper]{report}
\usepackage[utf8]{inputenc}
\usepackage[T5]{fontenc}
\usepackage[vietnamese]{babel}
\usepackage[a4paper,margin=2.5cm]{geometry}
\usepackage{longtable,booktabs,tabularx,array,multirow,xcolor}
\usepackage{enumitem}
\usepackage{ragged2e}
\usepackage{graphicx}
\usepackage{verbatim}
\usepackage{etoolbox}
\usepackage[hidelinks,unicode,linktoc=all]{hyperref}
\usepackage{bookmark}
\AtBeginEnvironment{longtable}{\tiny\setlength{\tabcolsep}{0pt}\hfuzz=1000pt}
\AtEndEnvironment{longtable}{\normalsize}
\ExplSyntaxOn
\makeatletter
\makeatother
\ExplSyntaxOff
\makeatletter
\g@addto@macro\UrlBreaks{\do\/\do\-\do\.\do\_\do\@\do\:}
\makeatother
\newcommand{\codepath}[1]{\begingroup\Urlmuskip=0mu plus 3mu\path{#1}\endgroup}
\setlength{\parindent}{0pt}
\setlength{\parskip}{6pt}
\setlength{\emergencystretch}{3em}
\newcolumntype{Y}{>{\raggedright\arraybackslash}X}
\setlength{\tabcolsep}{1.5pt}
\renewcommand{\arraystretch}{1.18}
\AtBeginEnvironment{longtable}{\scriptsize}
\AtEndEnvironment{longtable}{\normalsize}
\begin{document}
\pagenumbering{gobble}
% ==================== Báo Cáo Tổng Hợp Đồ Án ====================
\begin{titlepage}
\centering
\vspace*{1cm}
\rule{0.92\textwidth}{0.8pt}\par\vspace{2mm}\rule{0.90\textwidth}{0.4pt}
\vspace{0.8cm} \\
{\bfseries ĐẠI HỌC QUỐC GIA THÀNH PHỐ HỒ CHÍ MINH\par}
{\bfseries TRƯỜNG ĐẠI HỌC KHOA HỌC TỰ NHIÊN\par}
{\bfseries KHOA CÔNG NGHỆ THÔNG TIN\par}
\vspace{1cm}
% TODO: Chèn logo Trường Đại học Khoa học Tự nhiên tại đây
\includegraphics[width=3cm]{LogoKHTN.png} \\
\vspace{0.8cm}
{\Large\bfseries BÁO CÁO ĐỒ ÁN CUỐI KỲ\par}
\vspace{0.5cm}
{\Large\bfseries Đề tài: Xây dựng hệ thống thương mại điện tử bán voucher giảm giá trực tuyến\par}
\vfill
\begin{tabular}{rl}
Nhóm: & Nhóm đồ án Thương mại điện tử\\
Lớp: & CQ2023/1 \\
Giảng viên hướng dẫn: & Thầy Nguyễn Đức Huy\\
& Thầy Lương Vĩ Minh\\
Thành viên thực hiện: & 23120116 - Nguyễn Việt Cường\\
& 23120162 - Lê Hải Sơn\\
& 23120265 - Nguyễn Thái Hoàng\\
& 23120141 - Lê Văn Lộc
\end{tabular}
\vfill
{\bfseries TP. Hồ Chí Minh, năm 2026\par}
\vspace{0.8cm}
\rule{0.90\textwidth}{0.4pt}\par\vspace{2mm}\rule{0.92\textwidth}{0.8pt}
\end{titlepage}
\clearpage
\pagenumbering{roman}
\tableofcontents
\clearpage
\pagenumbering{arabic}

\chapter{Giới Thiệu Đề Tài}

\section{Giới thiệu}
Đề tài xây dựng hệ thống thương mại điện tử bán voucher giảm giá trực tuyến. Hệ thống cần hỗ trợ ba vai trò: khách hàng, đối tác và quản trị viên. Tài liệu BRD yêu cầu các luồng mua voucher, phát hành code, xác nhận sử dụng, báo cáo, quản lý nội dung và kiểm toán.

\section{Bối Cảnh Và Mục Tiêu}
Mục tiêu hệ thống là cung cấp kênh bán voucher trực tuyến có quy trình duyệt rõ ràng, kiểm soát tồn kho và trạng thái, có phân quyền và có khả năng báo cáo quản trị.

\section{Phạm Vi Hệ Thống}
\begin{itemize}[leftmargin=1.2cm]
  \item Khách hàng: tìm kiếm, giỏ hàng, mua, thanh toán, nhận voucher code, đánh giá và khiếu nại.
  \item Đối tác: đăng ký hồ sơ, tạo voucher, gửi duyệt, quản lý voucher, kiểm tra và xác nhận sử dụng.
  \item Quản trị viên: duyệt đối tác, duyệt voucher, quản lý đơn hàng, nội dung, dashboard và nhật ký.
\end{itemize}

\chapter{Tổng Quan Codebase}
\section{Công nghệ}
Frontend dùng React 18, React Router, Axios. Backend dùng Node.js, Express, PostgreSQL, JWT, bcryptjs. Bằng chứng: \codepath{frontend/package.json}, \codepath{backend/package.json}.

\section{Cấu trúc Repository}

Cấu trúc repository được tổ chức theo mô hình tách biệt frontend, backend và tài liệu báo cáo. Các thư mục sinh ra bởi quá trình cài đặt hoặc build như \texttt{node\_modules}, \texttt{dist}, \texttt{build}, \texttt{coverage}, cache và log không được liệt kê trong cây dưới đây.
\newpage
\begin{verbatim}
ec-voucher-system/
|-- backend/                         # REST API và xử lý nghiệp vụ phía server
|   |-- src/                         # Mã nguồn chính của backend
|   |   |-- app.js                   # Cấu hình Express app, middleware và routes
|   |   |-- index.js                 # Entry point khởi động server
|   |   |-- config/                  # Kết nối CSDL, schema và dữ liệu seed
|   |   |-- controllers/             # Xử lý request theo từng nghiệp vụ
|   |   |   `-- __tests__/           # Kiểm thử luồng nghiệp vụ backend
|   |   |-- middleware/              # Middleware xác thực, phân quyền, validate, lỗi
|   |   |-- models/                  # Các câu truy vấn SQL và lớp truy xuất dữ liệu
|   |   |-- routes/                  # Định nghĩa các REST API route
|   |   `-- utils/                   # Tiện ích dùng chung, tích hợp VNPay/PayPal
|   |-- TESTS/                       # Tài liệu/lệnh kiểm thử API thủ công
|   |-- MIDDLEWARE_GUIDE.md          # Ghi chú triển khai middleware
|   `-- package.json                 # Cấu hình npm và script backend
|-- frontend/                        # Ứng dụng React cho người dùng, đối tác, admin
|   |-- public/                      # Tài nguyên tĩnh public của frontend
|   |-- src/                         # Mã nguồn chính của frontend
|   |   |-- App.jsx                  # Cấu hình route và layout ứng dụng
|   |   |-- index.js                 # Entry point mount React app
|   |   |-- components/              # Component dùng lại như Navbar, ProtectedRoute
|   |   |-- context/                 # Global state cho xác thực và giỏ hàng
|   |   |-- hooks/                   # Custom hook cho logic giao diện
|   |   |-- pages/                   # Các trang Customer, Partner, Admin và trang chung
|   |   `-- services/                # Lớp gọi API backend bằng Axios
|   `-- package.json                 # Cấu hình npm và script frontend
|-- docs/                            # Tài liệu dự án và báo cáo
|   `-- latex/                       # Nguồn LaTeX của báo cáo tổng hợp
|       |-- revised/                 # Các chương báo cáo được tách thành file riêng
|       `-- 00_FULL_COMBINED_REPORT_REVISED.tex
|                                      # File báo cáo LaTeX tổng hợp
|-- .devcontainer/                   # Cấu hình môi trường phát triển bằng container
|-- docker-compose.yml               # Khởi chạy backend, frontend, database bằng Docker
|-- package.json                     # Cấu hình npm cấp root cho thao tác toàn dự án
|-- README.md                        # Hướng dẫn tổng quan dự án
|-- deployment_guide.md              # Hướng dẫn triển khai hệ thống
|-- Yeu_Cau_Nghiep_Vu.md             # Tài liệu yêu cầu nghiệp vụ chính
|-- yeu_cau_du_lieu_nghiep_vu.md     # Yêu cầu dữ liệu nghiệp vụ
|-- yeu_cau_phi_chuc_nang.md         # Yêu cầu phi chức năng
|-- quy_tac_nghiep_vu.md             # Quy tắc nghiệp vụ của hệ thống voucher
|-- TEST_CASES_PHASE2.md             # Danh sách test case giai đoạn 2
`-- test_accounts.md                 # Tài khoản mẫu phục vụ kiểm thử/demo
\end{verbatim}

\newpage
\subsection{Công nghệ Chi tiết}

\begin{itemize}[leftmargin=1.2cm]
  \item \textbf{Backend}: Node.js, Express.js, PostgreSQL (pg driver), JWT (jsonwebtoken), bcryptjs, UUID, Morgan (logging), CORS.
  \item \textbf{Frontend}: React 18, React Router DOM v6, Axios, react-icons, react-slider, CSS-in-JS.
  \item \textbf{Authentication}: JWT bearer tokens, bcryptjs password hashing, role-based authorization (ADMIN, PARTNER, CUSTOMER).
  \item \textbf{Payment Integration}: VNPay (IPN callback), PayPal SDK (Sandbox), thanh toán mô phỏng (COD).
  \item \textbf{Database}: PostgreSQL schema với enums (user\_role, voucher\_status, order\_status, issued\_status), constraints, indexes, foreign keys.
\end{itemize}

\chapter{Chức Năng Đã Xây Dựng}
Tóm tắt phạm vi đã triển khai (chi tiết use case, UI và đối chiếu yêu cầu xem các chương sau):
\begin{itemize}[leftmargin=1.2cm]
  \item \textbf{Xác thực}: đăng ký, đăng nhập, JWT, phân quyền ba vai trò.
  \item \textbf{Khách hàng}: tìm voucher, giỏ hàng, đặt hàng, thanh toán, nhận mã, đánh giá/khiếu nại.
  \item \textbf{Đối tác}: hồ sơ, chi nhánh, tạo/gửi duyệt voucher, quét mã, báo cáo.
  \item \textbf{Quản trị}: duyệt đối tác/voucher, đơn hàng, nội dung, dashboard, nhật ký (một phần).
\end{itemize}

\chapter{Mô Hình Nghiệp Vụ}
\section{Các Bên Liên Quan}
\begin{longtable}{>{\raggedright\arraybackslash}p{0.22\linewidth}>{\raggedright\arraybackslash}p{0.70\linewidth}}
\toprule
\textbf{Bên liên quan} & \textbf{Vai trò trong hệ thống}\\
\midrule
\endhead
Khách hàng (Customer) & Tìm kiếm, mua voucher, thanh toán, nhận mã, đánh giá, khiếu nại. Bằng chứng: \texttt{frontend/src/pages/VouchersPage.jsx}, \texttt{CartPage.jsx}.\\
Đối tác (Partner) & Đăng ký hồ sơ, tạo voucher, gửi duyệt, quét/xác nhận mã, xem báo cáo. Bằng chứng: \texttt{PartnerDashboardPage.jsx}, \texttt{partner.controller.js}.\\
Quản trị viên (Admin) & Duyệt đối tác/voucher, quản lý đơn, nội dung, khiếu nại, nhật ký. Bằng chứng: \texttt{AdminDashboardPage.jsx}, \texttt{admin.controller.js}.\\
Hệ thống thanh toán & VNPay/PayPal sandbox và thanh toán mô phỏng COD. Bằng chứng: \texttt{vnpay.js}, \texttt{paypal.js}, \texttt{order.controller.js}.\\
Cơ sở dữ liệu PostgreSQL & Lưu trữ thực thể nghiệp vụ. Bằng chứng: \texttt{backend/src/config/init.sql}.\\
\bottomrule
\end{longtable}

\section{Quy Trình Nghiệp Vụ Tổng Quát}
Luồng tổng quát: (1) Đối tác đăng ký và được admin duyệt; (2) Đối tác tạo voucher và gửi duyệt; (3) Admin duyệt voucher (RB-01); (4) Khách hàng tìm kiếm, mua, thanh toán; (5) Hệ thống phát hành mã (RB-05); (6) Khách dùng mã tại chi nhánh; (7) Đối tác xác nhận sử dụng; (8) Admin giám sát qua dashboard và nhật ký.

\textbf{NOTE: Nhóm cần bổ sung hình BPMN hoặc Activity Diagram cho luồng nghiệp vụ tại đây.}

\section{Luồng Nghiệp Vụ Mua Voucher}
Khách hàng duyệt danh sách voucher đã APPROVED (\texttt{listVouchers}), xem chi tiết, thêm vào giỏ (\texttt{CartContext}), tạo đơn (\texttt{createOrder}), chọn phương thức thanh toán và hoàn tất (\texttt{payOrder} hoặc VNPay/PayPal). Hệ thống kiểm RB-01, RB-04, RB-15 trước khi ghi nhận đơn.

\textbf{NOTE: Nhóm cần bổ sung Activity Diagram/BPMN dạng hình cho luồng này.}

\section{Luồng Đối Tác Tạo Và Gửi Duyệt Voucher}
Đối tác APPROVED truy cập \texttt{PartnerVoucherForm.jsx}, nhập thông tin (giá, stock, thời gian, chi nhánh), lưu DRAFT hoặc \texttt{PENDING\_APPROVAL}. Admin xem \texttt{/admin/vouchers} và phê duyệt.

\section{Luồng Quản Trị Viên Duyệt Voucher}
Admin lấy danh sách \texttt{GET /api/admin/vouchers/pending}, xem chi tiết, gọi \texttt{approveVoucher} hoặc \texttt{rejectVoucher}. Voucher APPROVED mới hiển thị cho khách hàng.

\section{Luồng Sử Dụng Và Xác Thực Voucher}
Khách nhận mã từ \texttt{MyVouchersPage}. Đối tác nhập mã tại \texttt{PartnerVoucherScan.jsx}, gọi \texttt{checkVoucher} (tra cứu) hoặc \texttt{scanVoucher} (xác nhận USED). RB-07, RB-08, RB-09 được kiểm tra tại \texttt{partner.controller.js}.

\section{Luồng Khiếu Nại}
Khách gửi khiếu nại qua \texttt{POST /api/users/complaints} (\texttt{user.controller.js}). Admin xử lý qua \texttt{GET/PATCH /api/admin/complaints}. Đối tác bị khóa có thể gửi \texttt{partner\_appeals}.

\section{Nhận Xét Liên Hệ Nghiệp Vụ Và Codebase}
Các trạng thái nghiệp vụ được ánh xạ rõ sang enum PostgreSQL (\texttt{voucher\_status}, \texttt{order\_status}, \texttt{issued\_voucher\_status}). Logic nghiệp vụ tập trung ở controller layer; frontend phản ánh qua \texttt{ProtectedRoute} và \texttt{usePartnerStatus}. Điểm còn thiếu so với BRD: job tự động hết hạn, audit log đầy đủ, refund tự động — \textbf{Chưa tìm thấy trong codebase} hoặc chỉ \textbf{một phần}.


\chapter{Đối Chiếu Yêu Cầu}
% Đối chiếu yêu cầu — tách riêng BR, BR-CUS, BR-PAR, BR-ADM, RB
\section{Tổng Quan Đối Chiếu}
Chương này đối chiếu yêu cầu nghiệp vụ (BR), yêu cầu theo vai trò (BR-CUS, BR-PAR, BR-ADM) và quy tắc nghiệp vụ (RB) với codebase hiện có. \textbf{Lưu ý:} BR là yêu cầu nghiệp vụ tổng thể; RB là quy tắc ràng buộc nghiệp vụ — hai nhóm mã \textbf{không} thay thế cho nhau.

\small
\subsection{Bảng Đối Chiếu Yêu Cầu Tổng Thể (BR-01 đến BR-07)}
\begin{longtable}{>{\raggedright\arraybackslash}p{0.08\linewidth}>{\raggedright\arraybackslash}p{0.18\linewidth}>{\raggedright\arraybackslash}p{0.10\linewidth}>{\raggedright\arraybackslash}p{0.28\linewidth}>{\raggedright\arraybackslash}p{0.14\linewidth}>{\raggedright\arraybackslash}p{0.14\linewidth}}
\toprule
\textbf{Mã} & \textbf{Nội dung} & \textbf{Tình trạng} & \textbf{Bằng chứng codebase} & \textbf{Nhận xét} & \textbf{Đề xuất}\\
\midrule
\endhead
BR-01 & Quản lý tài khoản người dùng & Hoàn thành & \texttt{auth.routes.js}, \texttt{auth.controller.js}, \texttt{user.routes.js}, \texttt{RegisterPage.jsx}, \texttt{LoginPage.jsx}, \texttt{ProfilePage.jsx}, bảng \texttt{users} & Đủ đăng ký, đăng nhập, đổi/quên mật khẩu, JWT, phân quyền & Bổ sung xác thực email thật nếu cần\\
BR-02 & Quản lý danh mục và nội dung voucher & Một phần & \texttt{voucher.controller.js}, \texttt{admin.controller.js}, bảng \texttt{categories}, \texttt{banners}, \texttt{content\_pages}, \texttt{popups} & Có CRUD nội dung admin; tạm ngưng voucher qua \texttt{updateVoucherStatus} & Kiểm tra đủ trạng thái SUSPENDED/EXPIRED tự động\\
BR-03 & Mua hàng trực tuyến & Hoàn thành & \texttt{CartContext.jsx}, \texttt{CartPage.jsx}, \texttt{order.controller.js}, \texttt{order.routes.js}, bảng \texttt{orders}, \texttt{order\_items} & Giỏ hàng localStorage, tạo đơn, COD/VNPay/PayPal & —\\
BR-04 & Phát hành và quản lý voucher code & Hoàn thành & \texttt{order.controller.js} (\texttt{generateCode}), \texttt{order.queries.js}, bảng \texttt{issued\_vouchers} & Sinh mã sau thanh toán, theo dõi UNUSED/USED/EXPIRED & Job tự động chuyển EXPIRED\\
BR-05 & Kiểm tra và xác thực voucher & Hoàn thành & \texttt{partner.controller.js} (\texttt{checkVoucher}, \texttt{scanVoucher}), \texttt{PartnerVoucherScan.jsx}, route POST \texttt{/partner/vouchers/check} & Kiểm tra và xác nhận sử dụng tại chi nhánh & Bổ sung quét QR thật nếu yêu cầu\\
BR-06 & Kiểm duyệt và giám sát hệ thống & Hoàn thành & \texttt{admin.controller.js}, \texttt{AdminDashboardPage.jsx}, \texttt{AdminVoucherReview.jsx}, \texttt{admin.routes.js} & Duyệt đối tác/voucher, quản lý đơn, khiếu nại & Tách module admin lớn\\
BR-07 & Báo cáo và phân tích & Một phần & \texttt{admin.controller.js} (\texttt{getDashboard}), \texttt{partner.controller.js} (\texttt{getReports}), \texttt{PartnerReports.jsx} & Dashboard admin/partner có số liệu tổng hợp & \textbf{NOTE: Nhóm cần bổ sung số liệu demo chính xác trên slide}\\
\bottomrule
\end{longtable}

\subsection{Bảng Đối Chiếu Yêu Cầu Khách Hàng (BR-CUS-01 đến BR-CUS-08)}
\begin{longtable}{>{\raggedright\arraybackslash}p{0.10\linewidth}>{\raggedright\arraybackslash}p{0.16\linewidth}>{\raggedright\arraybackslash}p{0.10\linewidth}>{\raggedright\arraybackslash}p{0.28\linewidth}>{\raggedright\arraybackslash}p{0.14\linewidth}>{\raggedright\arraybackslash}p{0.14\linewidth}}
\toprule
\textbf{Mã} & \textbf{Nội dung} & \textbf{Tình trạng} & \textbf{Bằng chứng codebase} & \textbf{Nhận xét} & \textbf{Đề xuất}\\
\midrule
\endhead
BR-CUS-01 & Đăng ký tài khoản & Hoàn thành & \texttt{POST /api/auth/register}, \texttt{auth.controller.js}, \texttt{RegisterPage.jsx}, \texttt{users} & Kiểm tra trùng email/phone & Test case trùng email\\
BR-CUS-02 & Đăng nhập và quản lý hồ sơ & Hoàn thành & \texttt{LoginPage.jsx}, \texttt{ProfilePage.jsx}, \texttt{ForgotPasswordPage.jsx}, \texttt{user.controller.js} & JWT + \texttt{AuthContext.jsx} & —\\
BR-CUS-03 & Tìm kiếm voucher & Hoàn thành & \texttt{VouchersPage.jsx}, \texttt{voucher.controller.js} (\texttt{listVouchers}), \texttt{voucher.queries.js} & Lọc giá, danh mục, từ khóa & Lọc khu vực: Cần kiểm tra thêm\\
BR-CUS-04 & Xem chi tiết voucher & Hoàn thành & \texttt{VoucherDetailPage.jsx}, \texttt{GET /api/vouchers/:id} & Hiển thị giá, hạn, chi nhánh, đánh giá & Chính sách hoàn hủy: một phần qua \texttt{terms}\\
BR-CUS-05 & Quản lý giỏ hàng & Hoàn thành & \texttt{CartContext.jsx}, \texttt{CartPage.jsx} & Thêm/sửa/xóa, tổng tiền localStorage & —\\
BR-CUS-06 & Tạo đơn hàng & Hoàn thành & \texttt{order.controller.js} (\texttt{createOrder}), \texttt{CartPage.jsx} & Người nhận quà, payment method & —\\
BR-CUS-07 & Nhận voucher đã mua & Hoàn thành & \texttt{MyVouchersPage.jsx}, \texttt{user.controller.js} (\texttt{getMyVouchers}), \texttt{issued\_vouchers} & Xem code sau PAID; QR mô phỏng trên UI & \textbf{NOTE: Nhóm cần chụp màn hình QR/code}\\
BR-CUS-08 & Đánh giá và phản hồi & Hoàn thành & \texttt{review.controller.js}, \texttt{user.routes.js} (complaints), \texttt{complaint.service.js} & Review gắn \texttt{issued\_voucher\_id}; khiếu nại qua \texttt{complaints} & —\\
\bottomrule
\end{longtable}

\subsection{Bảng Đối Chiếu Yêu Cầu Đối Tác (BR-PAR-01 đến BR-PAR-07)}
\begin{longtable}{>{\raggedright\arraybackslash}p{0.10\linewidth}>{\raggedright\arraybackslash}p{0.16\linewidth}>{\raggedright\arraybackslash}p{0.10\linewidth}>{\raggedright\arraybackslash}p{0.28\linewidth}>{\raggedright\arraybackslash}p{0.14\linewidth}>{\raggedright\arraybackslash}p{0.14\linewidth}}
\toprule
\textbf{Mã} & \textbf{Nội dung} & \textbf{Tình trạng} & \textbf{Bằng chứng codebase} & \textbf{Nhận xét} & \textbf{Đề xuất}\\
\midrule
\endhead
BR-PAR-01 & Đăng ký và quản lý hồ sơ đối tác & Hoàn thành & \texttt{partner.controller.js} (\texttt{registerPartner}, \texttt{updateProfile}), \texttt{PartnerDashboardPage.jsx}, \texttt{partners}, \texttt{partner\_branches} & Đăng ký kèm user PARTNER & —\\
BR-PAR-02 & Tạo voucher & Hoàn thành & \texttt{voucher.controller.js} (\texttt{createVoucher}), \texttt{PartnerVoucherForm.jsx} & Validate RB-02, RB-03 & —\\
BR-PAR-03 & Gửi duyệt voucher & Hoàn thành & \texttt{createVoucher}/\texttt{updateVoucher} status \texttt{PENDING\_APPROVAL}, \texttt{PartnerVouchers.jsx} & Admin duyệt qua \texttt{approveVoucher} & —\\
BR-PAR-04 & Quản lý voucher & Hoàn thành & \texttt{PartnerVouchers.jsx}, \texttt{PUT /api/vouchers/:id} & Sửa trước duyệt; xem stock, trạng thái & —\\
BR-PAR-05 & Kiểm tra voucher code & Hoàn thành & \texttt{checkVoucher}, \texttt{PartnerVoucherScan.jsx} & Nhập mã, xem trạng thái & Quét QR camera: Cần kiểm tra thêm\\
BR-PAR-06 & Xác nhận sử dụng voucher & Hoàn thành & \texttt{scanVoucher}, cập nhật \texttt{issued\_vouchers.status=USED} & RB-07, RB-09 trong controller & —\\
BR-PAR-07 & Báo cáo đối tác & Hoàn thành & \texttt{getReports}, \texttt{PartnerReports.jsx}, \texttt{partner.queries.js} & Doanh thu, bán, dùng theo voucher & —\\
\bottomrule
\end{longtable}

\subsection{Bảng Đối Chiếu Yêu Cầu Quản Trị Viên (BR-ADM-01 đến BR-ADM-07)}
\begin{longtable}{>{\raggedright\arraybackslash}p{0.10\linewidth}>{\raggedright\arraybackslash}p{0.16\linewidth}>{\raggedright\arraybackslash}p{0.10\linewidth}>{\raggedright\arraybackslash}p{0.28\linewidth}>{\raggedright\arraybackslash}p{0.14\linewidth}>{\raggedright\arraybackslash}p{0.14\linewidth}}
\toprule
\textbf{Mã} & \textbf{Nội dung} & \textbf{Tình trạng} & \textbf{Bằng chứng codebase} & \textbf{Nhận xét} & \textbf{Đề xuất}\\
\midrule
\endhead
BR-ADM-01 & Quản lý người dùng & Hoàn thành & \texttt{getUsers}, \texttt{updateUserStatus}, \texttt{updateUserRole}, \texttt{AdminDashboardPage.jsx} & Khóa/mở, đổi role & —\\
BR-ADM-02 & Quản lý đối tác & Hoàn thành & \texttt{approvePartner}, \texttt{rejectPartner}, \texttt{updatePartnerStatusAny}, \texttt{getPartnerBranches} & Duyệt PENDING, SUSPENDED & —\\
BR-ADM-03 & Duyệt voucher & Hoàn thành & \texttt{approveVoucher}, \texttt{rejectVoucher}, \texttt{AdminVoucherReview.jsx} & RB-01 khi APPROVED mới bán & —\\
BR-ADM-04 & Quản lý đơn hàng & Một phần & \texttt{getOrders}, \texttt{updateOrderStatus}, bảng \texttt{orders} & Có REFUNDED enum; hoàn tiền mô phỏng thủ công & Hoàn tiền tự động: Chưa có\\
BR-ADM-05 & Quản lý nội dung & Hoàn thành & CRUD \texttt{categories/banners/pages/popups} trong \texttt{admin.controller.js} & Public qua \texttt{content.service.js} & —\\
BR-ADM-06 & Dashboard quản trị & Hoàn thành & \texttt{getDashboard}, \texttt{AdminDashboardPage.jsx} & Tổng hợp user, order, voucher & —\\
BR-ADM-07 & Nhật ký hệ thống & Một phần & \texttt{getLogs}, bảng \texttt{system\_logs}, ghi log tại \texttt{partner.controller.js}, \texttt{admin.queries.js} & Chưa phủ mọi thao tác admin & Mở rộng \texttt{INSERT INTO system\_logs}\\
\bottomrule
\end{longtable}

\subsection{Bảng Đối Chiếu Quy Tắc Nghiệp Vụ (RB-01 đến RB-15)}
\begin{longtable}{>{\raggedright\arraybackslash}p{0.08\linewidth}>{\raggedright\arraybackslash}p{0.20\linewidth}>{\raggedright\arraybackslash}p{0.10\linewidth}>{\raggedright\arraybackslash}p{0.26\linewidth}>{\raggedright\arraybackslash}p{0.14\linewidth}>{\raggedright\arraybackslash}p{0.14\linewidth}}
\toprule
\textbf{Mã} & \textbf{Nội dung} & \textbf{Tình trạng} & \textbf{Bằng chứng codebase} & \textbf{Nhận xét} & \textbf{Đề xuất}\\
\midrule
\endhead
RB-01 & Voucher chỉ bán khi đã duyệt & Hoàn thành & \texttt{voucher.controller.js} lọc APPROVED; \texttt{order.controller.js} kiểm \texttt{status !== APPROVED} & Đúng mapping RB, không nhầm BR & Test TC mua voucher chưa duyệt\\
RB-02 & Giá bán $<$ giá gốc & Hoàn thành & \texttt{chk\_price} trong \texttt{init.sql}; validate trong \texttt{voucher.controller.js} & DB + API & —\\
RB-03 & Thời gian bán và sử dụng rõ ràng & Một phần & Cột \texttt{sale\_start}, \texttt{sale\_end}, \texttt{valid\_until}; validate RB-03 & Kiểm tra lúc mua: Cần kiểm tra thêm & Bổ sung test hết hạn bán\\
RB-04 & Không bán khi hết SL/hết thời gian bán & Hoàn thành & \texttt{order.controller.js} kiểm stock và thời gian & \texttt{stock < quantity} → 409 & —\\
RB-05 & Code chỉ phát hành sau thanh toán & Hoàn thành & \texttt{payOrder}, \texttt{verifyVnpayPayment}; insert \texttt{issued\_vouchers} khi PAID & Comment RB-05 trong \texttt{init.sql} & —\\
RB-06 & Code duy nhất, khó đoán & Hoàn thành & \texttt{crypto.randomBytes(8)} trong \texttt{order.controller.js}; UNIQUE \texttt{code} & Prefix VCH- & —\\
RB-07 & Voucher đã dùng không dùng lại & Hoàn thành & \texttt{scanVoucher} kiểm \texttt{UNUSED}; test trong \texttt{business-flows.test.js} & — & —\\
RB-08 & Hết hạn/hủy/khóa không dùng & Một phần & Kiểm tra \texttt{expires\_at}, status trong \texttt{scanVoucher} & Job EXPIRED tự động: Chưa tìm thấy trong codebase & Cron job hết hạn\\
RB-09 & Đối tác chỉ xác thực voucher của mình & Hoàn thành & \texttt{checkVoucher}/\texttt{scanVoucher} so \texttt{partner\_id} & Test role partner khác & —\\
RB-10 & Chỉ đánh giá sau khi mua/dùng & Hoàn thành & \texttt{review.controller.js} kiểm \texttt{issued\_voucher} thuộc customer & UNIQUE \texttt{issued\_voucher\_id} & —\\
RB-11 & Số bán không vượt phát hành & Hoàn thành & \texttt{updateVoucherStockQuery}, \texttt{chk\_stock} & Giảm stock khi thanh toán & —\\
RB-12 & Thao tác quản trị ghi nhật ký & Một phần & \texttt{system\_logs}, \texttt{getLogs}; log khi scan, một số admin action & Nhiều PATCH admin chưa log & Ghi log đồng nhất\\
RB-13 & Đơn hủy không phát hành voucher & Hoàn thành & Chỉ sinh code trong nhánh PAID; \texttt{order\_status} CANCELLED & — & Test đơn PENDING\\
RB-14 & Chính sách hoàn tiền theo voucher/sàn & Một phần & \texttt{terms} text; \texttt{REFUNDED} status thủ công & Tự động hóa refund: Chưa có & Mô phỏng sâu hơn nếu cần\\
RB-15 & Kiểm tra tồn kho lúc đặt mua/thanh toán & Hoàn thành & \texttt{createOrder} và \texttt{payOrder} kiểm stock & Comment RB-15 trong schema & —\\
\bottomrule
\end{longtable}
\normalsize


\chapter{Cơ Sở Dữ Liệu Và Dữ Liệu Mẫu}
Schema: \codepath{backend/src/config/init.sql}; seed mở rộng: \texttt{seed-data.sql}. Các bảng chính liệt kê ở Chương ERD và Từ Điển Dữ Liệu. Dữ liệu mẫu gồm tài khoản admin, customer, partner và voucher demo để minh họa các tính năng chính của hệ thống.

\chapter{Đặc Tả Yêu Cầu Phần Mềm (SRS)}
\section{Giới Thiệu}
Tài liệu SRS mô tả yêu cầu phần mềm, ma trận truy vết và ràng buộc kỹ thuật. Chi tiết đối chiếu từng mã BR/RB xem Chương Đối Chiếu Yêu Cầu; luồng nghiệp vụ xem Chương Mô Hình Nghiệp Vụ và Đặc Tả Use Case.

\section{Mục Đích Và Phạm Vi}
Chuẩn hóa yêu cầu chức năng, phi chức năng, dữ liệu và truy vết tới module code. Phạm vi: web app React + Express API + PostgreSQL cho ba vai trò CUSTOMER, PARTNER, ADMIN.

\section{Thuật Ngữ}
\begin{longtable}{>{\raggedright\arraybackslash}p{0.22\linewidth}>{\raggedright\arraybackslash}p{0.72\linewidth}}
\toprule
\textbf{Thuật ngữ} & \textbf{Mô tả}\\
\midrule
BR & Yêu cầu nghiệp vụ (Business Requirement).\\
RB & Quy tắc nghiệp vụ (Business Rule).\\
Issued voucher & Mã voucher điện tử phát hành sau thanh toán thành công.\\
\bottomrule
\end{longtable}

\section{Yêu Cầu Phi Chức Năng}
\begin{itemize}[leftmargin=1.2cm]
  \item \textbf{Bảo mật}: JWT (\texttt{auth.middleware.js}), bcrypt (\texttt{auth.controller.js}), phân quyền \texttt{authorize()}.
  \item \textbf{Hiệu năng}: Index trên \texttt{vouchers}, \texttt{orders}, \texttt{issued\_vouchers} trong \texttt{init.sql}.
  \item \textbf{Kiểm toán}: Bảng \texttt{system\_logs} — mức độ phủ \textbf{một phần} (xem BR-ADM-07, RB-12).
  \item \textbf{Triển khai}: Docker Compose (xem Chương Hướng Dẫn Cài Đặt).
\end{itemize}

\section{Yêu Cầu Dữ Liệu Tóm Tắt}
Thực thể chính: \texttt{users}, \texttt{partners}, \texttt{vouchers}, \texttt{orders}, \texttt{issued\_vouchers}, \texttt{reviews}, \texttt{complaints}. Chi tiết ERD và từ điển dữ liệu ở các chương riêng.

\section{Ràng Buộc Kỹ Thuật}
RB-02 (giá), RB-05 (phát hành sau thanh toán), RB-11/RB-15 (tồn kho) được enforce tại API và/hoặc constraint DB. Frontend không được tin tưởng dữ liệu giỏ hàng — backend tái kiểm tra khi \texttt{createOrder}/\texttt{payOrder}.

\section{Tiêu Chí Nghiệm Thu}
Ba vai trò hoạt động; dữ liệu seed; luồng mua--phát hành--sử dụng; tài liệu đồ án (báo cáo, SRS, use case, ERD, test plan).

\section{Ma Trận Truy Vết (Tóm Tắt)}
Ma trận đầy đủ nằm ở chương Đối Chiếu Yêu Cầu. Bảng rút gọn dưới đây liên kết use case quan trọng tới module.

\small
\begin{longtable}{>{\raggedright\arraybackslash}p{0.12\linewidth}>{\raggedright\arraybackslash}p{0.14\linewidth}>{\raggedright\arraybackslash}p{0.22\linewidth}>{\raggedright\arraybackslash}p{0.34\linewidth}>{\raggedright\arraybackslash}p{0.10\linewidth}}
\toprule
\textbf{BR/RB} & \textbf{UC} & \textbf{Module} & \textbf{File chính} & \textbf{TT}\\
\midrule
\endhead
BR-01 & UC-CUS-02 & Auth & \texttt{auth.controller.js}, \texttt{users} & HT\\
BR-03 & UC-CUS-06 & Order & \texttt{order.controller.js}, \texttt{CartPage.jsx} & HT\\
BR-04 & UC-CUS-07 & Issued & \texttt{order.controller.js}, \texttt{issued\_vouchers} & HT\\
RB-01 & UC-CUS-06 & Voucher & \texttt{voucher.controller.js}, \texttt{order.controller.js} & HT\\
RB-05 & UC-CUS-07 & Payment & \texttt{payOrder}, \texttt{issued\_vouchers} & HT\\
BR-PAR-02 & UC-PAR-02 & Partner Voucher & \texttt{PartnerVoucherForm.jsx} & HT\\
BR-ADM-03 & UC-ADM-02 & Admin & \texttt{AdminVoucherReview.jsx} & HT\\
BR-ADM-07 & UC-ADM-05 & Logs & \texttt{getLogs}, \texttt{system\_logs} & MP\\
\bottomrule
\end{longtable}
\normalsize
\textit{Chú thích TT: HT = Hoàn thành; MP = Một phần.}


\chapter{Đặc Tả Use Case}
\section{Tổng Quan}
Hệ thống có ba actor: Khách hàng, Đối tác, Quản trị viên. Bảng tổng hợp liệt kê use case chính; mỗi use case quan trọng có đặc tả chi tiết theo template thống nhất.

\section{Bảng Tổng Hợp Use Case}
\small
\begin{longtable}{>{\raggedright\arraybackslash}p{0.14\linewidth}>{\raggedright\arraybackslash}p{0.26\linewidth}>{\raggedright\arraybackslash}p{0.14\linewidth}>{\raggedright\arraybackslash}p{0.34\linewidth}}
\toprule
\textbf{Mã} & \textbf{Tên} & \textbf{Actor} & \textbf{Bằng chứng codebase}\\
\midrule
\endhead
UC-CUS-01 & Đăng ký tài khoản & Khách hàng & \texttt{auth.routes.js}, \texttt{RegisterPage.jsx}\\
UC-CUS-02 & Đăng nhập & Khách hàng & \texttt{LoginPage.jsx}, \texttt{auth.controller.js}\\
UC-CUS-03 & Tìm kiếm/lọc voucher & Khách hàng & \texttt{VouchersPage.jsx}, \texttt{listVouchers}\\
UC-CUS-04 & Xem chi tiết voucher & Khách hàng & \texttt{VoucherDetailPage.jsx}\\
UC-CUS-05 & Thêm vào giỏ hàng & Khách hàng & \texttt{CartContext.jsx}\\
UC-CUS-06 & Tạo đơn hàng & Khách hàng & \texttt{order.controller.js}, \texttt{CartPage.jsx}\\
UC-CUS-07 & Thanh toán và nhận code & Khách hàng & \texttt{payOrder}, \texttt{MyVouchersPage.jsx}\\
UC-CUS-08 & Đánh giá/khiếu nại & Khách hàng & \texttt{review.controller.js}, \texttt{user.routes.js}\\
UC-CUS-09 & Quản lý hồ sơ cá nhân & Khách hàng & \texttt{ProfilePage.jsx}, \texttt{user.routes.js}, \texttt{auth.routes.js}\\
UC-CUS-10 & Xem đơn hàng/voucher đã mua & Khách hàng & \texttt{OrdersPage.jsx}, \texttt{MyVouchersPage.jsx}, \texttt{listMyOrders}\\
UC-CUS-11 & Xem nội dung công khai & Khách hàng & \texttt{ContentPage.jsx}, \texttt{content.service.js}\\
UC-PAR-01 & Đăng ký hồ sơ đối tác & Đối tác & \texttt{registerPartner}, \texttt{RegisterPage.jsx}\\
UC-PAR-02 & Tạo voucher & Đối tác & \texttt{createVoucher}, \texttt{PartnerVoucherForm.jsx}\\
UC-PAR-03 & Gửi duyệt voucher & Đối tác & Status \texttt{PENDING\_APPROVAL}\\
UC-PAR-04 & Kiểm tra voucher code & Đối tác & \texttt{checkVoucher}, \texttt{PartnerVoucherScan.jsx}\\
UC-PAR-05 & Xác nhận sử dụng & Đối tác & \texttt{scanVoucher}\\
UC-PAR-06 & Quản lý hồ sơ và chi nhánh & Đối tác & \texttt{updatePartnerProfile}, \texttt{createBranch}, \texttt{updatePartnerBranch}\\
UC-PAR-07 & Xem dashboard/báo cáo & Đối tác & \texttt{getPartnerDashboard}, \texttt{getPartnerReport}, \texttt{PartnerReports.jsx}\\
UC-PAR-08 & Gửi khiếu nại mở khóa & Đối tác & \texttt{createPartnerAppeal}, \texttt{PartnerDashboardPage.jsx}\\
UC-PAR-09 & Phản hồi đánh giá & Đối tác & \texttt{replyReview}, \texttt{review.routes.js}\\
UC-ADM-01 & Duyệt đối tác & Admin & \texttt{approvePartner}, \texttt{AdminDashboardPage.jsx}\\
UC-ADM-02 & Duyệt/từ chối voucher & Admin & \texttt{approveVoucher}, \texttt{AdminVoucherReview.jsx}\\
UC-ADM-03 & Quản lý đơn hàng & Admin & \texttt{getOrders}, \texttt{updateOrderStatus}\\
UC-ADM-04 & Xử lý khiếu nại & Admin & \texttt{updateComplaintStatus}\\
UC-ADM-05 & Dashboard/nhật ký & Admin & \texttt{getDashboard}, \texttt{getLogs}\\
UC-ADM-06 & Quản lý người dùng & Admin & \texttt{getUsers}, \texttt{updateUserStatus}, \texttt{updateUserRole}\\
UC-ADM-07 & Quản lý nội dung & Admin & CRUD \texttt{categories/banners/pages/popups}\\
UC-ADM-08 & Quản lý đối tác mở rộng & Admin & \texttt{updatePartnerStatusAny}, \texttt{getPartnerBranches}, \texttt{updatePartnerAppealStatus}\\
\bottomrule
\end{longtable}
\normalsize

% Macro-like repeated structure for use cases
\newcommand{\uctitle}[2]{\subsection{#1 --- #2}}

\uctitle{UC-CUS-01}{Đăng Ký Tài Khoản}
\begin{longtable}{>{\raggedright\arraybackslash}p{0.28\linewidth}>{\raggedright\arraybackslash}p{0.66\linewidth}}
Actor chính & Khách hàng (hoặc Đối tác khi chọn role PARTNER).\\
Mục tiêu & Tạo tài khoản mới trên hệ thống.\\
Tiền điều kiện & Email/phone chưa tồn tại.\\
Hậu điều kiện & Bản ghi \texttt{users} mới; JWT trả về (nếu auto-login).\\
Luồng chính & 1. Mở \texttt{/register}. 2. Nhập thông tin. 3. POST \texttt{/api/auth/register}. 4. Hệ thống hash mật khẩu, insert user. 5. Hiển thị thành công.\\
Luồng thay thế & A1: Đăng ký đối tác → tạo thêm \texttt{partners} PENDING.\\
Ngoại lệ & E1: Trùng email/phone → lỗi validation.\\
Dữ liệu đầu vào & email/phone, password, full\_name, role.\\
Dữ liệu đầu ra & user object, token.\\
RB liên quan & —\\
Bằng chứng & \texttt{auth.controller.js} (\texttt{register}), \texttt{RegisterPage.jsx}, bảng \texttt{users}.\\
Ghi chú & Xác thực email mô phỏng.\\
\end{longtable}

\uctitle{UC-CUS-02}{Đăng Nhập}
\begin{longtable}{>{\raggedright\arraybackslash}p{0.28\linewidth}>{\raggedright\arraybackslash}p{0.66\linewidth}}
Actor chính & Khách hàng / Đối tác / Admin.\\
Mục tiêu & Xác thực và nhận phiên JWT.\\
Tiền điều kiện & Tài khoản tồn tại, \texttt{is\_active=true}.\\
Hậu điều kiện & Token lưu localStorage qua \texttt{AuthContext}.\\
Luồng chính & 1. POST \texttt{/api/auth/login}. 2. So khớp bcrypt. 3. Trả JWT. 4. Redirect theo role.\\
Ngoại lệ & E1: Sai mật khẩu → 401. E2: Tài khoản khóa.\\
RB liên quan & —\\
Bằng chứng & \texttt{LoginPage.jsx}, \texttt{auth.controller.js}, \texttt{ProtectedRoute.jsx}.\\
\end{longtable}

\uctitle{UC-CUS-03}{Tìm Kiếm Và Lọc Voucher}
\begin{longtable}{>{\raggedright\arraybackslash}p{0.28\linewidth}>{\raggedright\arraybackslash}p{0.66\linewidth}}
Actor chính & Khách hàng.\\
Mục tiêu & Tìm voucher phù hợp.\\
Tiền điều kiện & Có voucher APPROVED.\\
Luồng chính & 1. \texttt{VouchersPage} gọi GET \texttt{/api/vouchers}. 2. Truyền query lọc/sort. 3. Hiển thị danh sách.\\
Ngoại lệ & E1: Không có kết quả.\\
RB liên quan & RB-01 (chỉ APPROVED).\\
Bằng chứng & \texttt{voucher.controller.js} (\texttt{listVouchers}), \texttt{VouchersPage.jsx}.\\
Ghi chú & Lọc khu vực: \textbf{Cần kiểm tra thêm}.\\
\end{longtable}

\uctitle{UC-CUS-04}{Xem Chi Tiết Voucher}
\begin{longtable}{>{\raggedright\arraybackslash}p{0.28\linewidth}>{\raggedright\arraybackslash}p{0.66\linewidth}}
Bằng chứng & \texttt{VoucherDetailPage.jsx}, \texttt{getVoucherById}, \texttt{review.routes.js}.\\
RB liên quan & RB-01.\\
Luồng chính & Hiển thị giá, hạn, chi nhánh (\texttt{voucher\_applicable\_branches}), đánh giá.\\
\end{longtable}

\uctitle{UC-CUS-05}{Thêm Voucher Vào Giỏ Hàng}
\begin{longtable}{>{\raggedright\arraybackslash}p{0.28\linewidth}>{\raggedright\arraybackslash}p{0.66\linewidth}}
Bằng chứng & \texttt{CartContext.jsx} (\texttt{addItem}), \texttt{VoucherDetailPage.jsx}.\\
Hậu điều kiện & Item trong localStorage giỏ hàng.\\
Ngoại lệ & E1: Chưa đăng nhập CUSTOMER khi checkout (chặn ở route).\\
\end{longtable}

\uctitle{UC-CUS-06}{Tạo Đơn Hàng}
\begin{longtable}{>{\raggedright\arraybackslash}p{0.28\linewidth}>{\raggedright\arraybackslash}p{0.66\linewidth}}
Luồng chính & POST \texttt{/api/orders} với items từ giỏ; kiểm APPROVED, stock, thời gian bán.\\
Hậu điều kiện & Order \texttt{PENDING}.\\
RB liên quan & RB-01, RB-04, RB-15.\\
Bằng chứng & \texttt{order.controller.js} (\texttt{createOrder}), \texttt{CartPage.jsx}.\\
\end{longtable}

\uctitle{UC-CUS-07}{Thanh Toán Và Nhận Voucher Code}
\begin{longtable}{>{\raggedright\arraybackslash}p{0.28\linewidth}>{\raggedright\arraybackslash}p{0.66\linewidth}}
Luồng chính & COD: \texttt{payOrder}; VNPay: \texttt{createVnpayPaymentUrl} → callback; sau PAID sinh \texttt{issued\_vouchers}.\\
RB liên quan & RB-05, RB-06, RB-11.\\
Bằng chứng & \texttt{order.controller.js}, \texttt{MyVouchersPage.jsx}, \texttt{VnpayReturnPage.jsx}.\\
\end{longtable}

\uctitle{UC-CUS-08}{Gửi Đánh Giá/Khiếu Nại}
\begin{longtable}{>{\raggedright\arraybackslash}p{0.28\linewidth}>{\raggedright\arraybackslash}p{0.66\linewidth}}
Luồng chính & Review: POST \texttt{/api/reviews}. Khiếu nại: POST \texttt{/api/users/complaints}.\\
RB liên quan & RB-10.\\
Bằng chứng & \texttt{review.controller.js}, \texttt{user.controller.js} (\texttt{createComplaint}).\\
\end{longtable}

\uctitle{UC-CUS-09}{Quản Lý Hồ Sơ Cá Nhân}
\begin{longtable}{>{\raggedright\arraybackslash}p{0.28\linewidth}>{\raggedright\arraybackslash}p{0.66\linewidth}}
Actor chính & Người dùng đã đăng nhập.\\
Mục tiêu & Xem/cập nhật thông tin cá nhân và đổi mật khẩu.\\
Tiền điều kiện & Có JWT hợp lệ.\\
Luồng chính & 1. Mở \texttt{/profile}. 2. GET \texttt{/api/users/profile}. 3. PUT \texttt{/api/users/profile} để cập nhật hồ sơ hoặc POST \texttt{/api/auth/change-password} để đổi mật khẩu.\\
Ngoại lệ & E1: Token hết hạn hoặc mật khẩu hiện tại sai.\\
Bằng chứng & \texttt{ProfilePage.jsx}, \texttt{user.controller.js}, \texttt{auth.controller.js}.\\
\end{longtable}

\uctitle{UC-CUS-10}{Xem Đơn Hàng Và Voucher Đã Mua}
\begin{longtable}{>{\raggedright\arraybackslash}p{0.28\linewidth}>{\raggedright\arraybackslash}p{0.66\linewidth}}
Actor chính & Khách hàng.\\
Mục tiêu & Theo dõi lịch sử đơn hàng và danh sách mã voucher đã phát hành.\\
Tiền điều kiện & Khách hàng đã đăng nhập.\\
Luồng chính & 1. Mở \texttt{/orders} hoặc \texttt{/my-vouchers}. 2. GET \texttt{/api/orders/my} hoặc GET \texttt{/api/users/vouchers}. 3. Hiển thị trạng thái đơn, item, code và trạng thái sử dụng.\\
Hậu điều kiện & Khách hàng nắm được mã voucher UNUSED/USED/EXPIRED.\\
Bằng chứng & \texttt{OrdersPage.jsx}, \texttt{MyVouchersPage.jsx}, \texttt{order.controller.js}, \texttt{user.controller.js}.\\
\end{longtable}

\uctitle{UC-CUS-11}{Xem Nội Dung Công Khai}
\begin{longtable}{>{\raggedright\arraybackslash}p{0.28\linewidth}>{\raggedright\arraybackslash}p{0.66\linewidth}}
Actor chính & Khách hàng hoặc người dùng chưa đăng nhập.\\
Mục tiêu & Xem banner, popup và các trang nội dung như điều khoản, chính sách, hướng dẫn.\\
Luồng chính & 1. Trang chủ lấy banner/popup đang active. 2. Người dùng mở \texttt{/pages/:slug}. 3. GET \texttt{/api/admin/content/pages/public/:slug}.\\
Bằng chứng & \texttt{HomePage.jsx}, \texttt{ContentPage.jsx}, \texttt{content.service.js}, \texttt{admin.routes.js}.\\
\end{longtable}

\uctitle{UC-PAR-01}{Đối Tác Đăng Ký Hồ Sơ}
\begin{longtable}{>{\raggedright\arraybackslash}p{0.28\linewidth}>{\raggedright\arraybackslash}p{0.66\linewidth}}
Hậu điều kiện & \texttt{partners.status=PENDING}.\\
Bằng chứng & \texttt{registerPartner}, \texttt{RegisterPage.jsx}, \texttt{PartnerDashboardPage.jsx}.\\
\end{longtable}

\uctitle{UC-PAR-02}{Đối Tác Tạo Voucher}
\begin{longtable}{>{\raggedright\arraybackslash}p{0.28\linewidth}>{\raggedright\arraybackslash}p{0.66\linewidth}}
Tiền điều kiện & Partner APPROVED (\texttt{usePartnerStatus.js}).\\
RB liên quan & RB-02, RB-03.\\
Bằng chứng & \texttt{createVoucher}, \texttt{PartnerVoucherForm.jsx}.\\
\end{longtable}

\uctitle{UC-PAR-03}{Đối Tác Gửi Duyệt Voucher}
\begin{longtable}{>{\raggedright\arraybackslash}p{0.28\linewidth}>{\raggedright\arraybackslash}p{0.66\linewidth}}
Luồng chính & Lưu voucher với \texttt{status=PENDING\_APPROVAL}.\\
Bằng chứng & \texttt{voucher.controller.js}, \texttt{PartnerVouchers.jsx}.\\
\end{longtable}

\uctitle{UC-PAR-04}{Đối Tác Kiểm Tra Voucher Code}
\begin{longtable}{>{\raggedright\arraybackslash}p{0.28\linewidth}>{\raggedright\arraybackslash}p{0.66\linewidth}}
Luồng chính & POST \texttt{/api/partner/vouchers/check} với code + branch\_id.\\
RB liên quan & RB-08, RB-09.\\
Bằng chứng & \texttt{checkVoucher}, \texttt{PartnerVoucherScan.jsx}.\\
\end{longtable}

\uctitle{UC-PAR-05}{Đối Tác Xác Nhận Sử Dụng Voucher}
\begin{longtable}{>{\raggedright\arraybackslash}p{0.28\linewidth}>{\raggedright\arraybackslash}p{0.66\linewidth}}
Luồng chính & POST scan → cập nhật USED, \texttt{used\_at}, \texttt{used\_at\_branch}.\\
RB liên quan & RB-07, RB-09.\\
Bằng chứng & \texttt{scanVoucher}, ghi \texttt{system\_logs}.\\
\end{longtable}

\uctitle{UC-PAR-06}{Quản Lý Hồ Sơ Và Chi Nhánh}
\begin{longtable}{>{\raggedright\arraybackslash}p{0.28\linewidth}>{\raggedright\arraybackslash}p{0.66\linewidth}}
Actor chính & Đối tác.\\
Mục tiêu & Cập nhật thông tin doanh nghiệp và quản lý chi nhánh áp dụng voucher.\\
Tiền điều kiện & Đối tác đã đăng nhập; một số thao tác chi nhánh yêu cầu status APPROVED.\\
Luồng chính & 1. Mở Partner Dashboard. 2. PUT \texttt{/api/partner/profile}. 3. POST \texttt{/api/partner/branches} hoặc PATCH \texttt{/api/partner/branches/:id}.\\
Ngoại lệ & E1: Partner chưa APPROVED bị chặn khi tạo/cập nhật chi nhánh.\\
Bằng chứng & \texttt{PartnerDashboardPage.jsx}, \texttt{partner.controller.js}, \texttt{partner.service.js}.\\
\end{longtable}

\uctitle{UC-PAR-07}{Xem Dashboard Và Báo Cáo}
\begin{longtable}{>{\raggedright\arraybackslash}p{0.28\linewidth}>{\raggedright\arraybackslash}p{0.66\linewidth}}
Actor chính & Đối tác.\\
Mục tiêu & Theo dõi voucher, doanh thu, số lượng bán/phát hành/sử dụng.\\
Luồng chính & 1. GET \texttt{/api/partner/dashboard}. 2. GET \texttt{/api/partner/reports}. 3. Hiển thị KPI và bảng báo cáo theo voucher.\\
Bằng chứng & \texttt{PartnerDashboardPage.jsx}, \texttt{PartnerReports.jsx}, \texttt{getPartnerDashboard}, \texttt{getPartnerReport}.\\
\end{longtable}

\uctitle{UC-PAR-08}{Gửi Khiếu Nại Mở Khóa}
\begin{longtable}{>{\raggedright\arraybackslash}p{0.28\linewidth}>{\raggedright\arraybackslash}p{0.66\linewidth}}
Actor chính & Đối tác bị tạm khóa.\\
Mục tiêu & Gửi yêu cầu admin xem xét mở khóa tài khoản đối tác.\\
Tiền điều kiện & Partner có status \texttt{SUSPENDED}.\\
Luồng chính & 1. Đối tác nhập tiêu đề, nội dung, bằng chứng. 2. POST \texttt{/api/partner/appeals}. 3. Hệ thống lưu \texttt{partner\_appeals.status=PENDING}.\\
Ngoại lệ & E1: Đối tác không bị khóa hoặc đã có appeal PENDING.\\
Bằng chứng & \texttt{createPartnerAppeal}, \texttt{PartnerDashboardPage.jsx}, bảng \texttt{partner\_appeals}.\\
\end{longtable}

\uctitle{UC-PAR-09}{Phản Hồi Đánh Giá}
\begin{longtable}{>{\raggedright\arraybackslash}p{0.28\linewidth}>{\raggedright\arraybackslash}p{0.66\linewidth}}
Actor chính & Đối tác.\\
Mục tiêu & Phản hồi review của khách hàng đối với voucher thuộc đối tác.\\
Tiền điều kiện & Partner APPROVED và review thuộc voucher của đối tác.\\
Luồng chính & 1. Đối tác nhập phản hồi. 2. POST \texttt{/api/reviews/:id/reply}. 3. Cập nhật \texttt{partner\_reply}.\\
Ngoại lệ & E1: Review không tồn tại hoặc không thuộc đối tác.\\
Bằng chứng & \texttt{replyReview}, \texttt{review.routes.js}, \texttt{review.controller.js}.\\
\end{longtable}

\uctitle{UC-ADM-01}{Quản Trị Viên Duyệt Đối Tác}
\begin{longtable}{>{\raggedright\arraybackslash}p{0.28\linewidth}>{\raggedright\arraybackslash}p{0.66\linewidth}}
Luồng chính & PATCH \texttt{/api/admin/partners/:id/approve|reject}.\\
Bằng chứng & \texttt{admin.controller.js}, tab Đối tác trong \texttt{AdminDashboardPage.jsx}.\\
\end{longtable}

\uctitle{UC-ADM-02}{Quản Trị Viên Duyệt/Từ Chối Voucher}
\begin{longtable}{>{\raggedright\arraybackslash}p{0.28\linewidth}>{\raggedright\arraybackslash}p{0.66\linewidth}}
RB liên quan & RB-01 (kích hoạt bán).\\
Bằng chứng & \texttt{approveVoucher}, \texttt{rejectVoucher}, \texttt{AdminVoucherReview.jsx}.\\
\end{longtable}

\uctitle{UC-ADM-03}{Quản Trị Viên Quản Lý Đơn Hàng}
\begin{longtable}{>{\raggedright\arraybackslash}p{0.28\linewidth}>{\raggedright\arraybackslash}p{0.66\linewidth}}
Bằng chứng & \texttt{getOrders}, \texttt{updateOrderStatus}, bảng \texttt{orders}.\\
Ghi chú & Hoàn tiền tự động: \textbf{Một phần} (REFUNDED thủ công).\\
\end{longtable}

\uctitle{UC-ADM-04}{Quản Trị Viên Xử Lý Khiếu Nại}
\begin{longtable}{>{\raggedright\arraybackslash}p{0.28\linewidth}>{\raggedright\arraybackslash}p{0.66\linewidth}}
Bằng chứng & \texttt{getComplaints}, \texttt{updateComplaintStatus}, bảng \texttt{complaints}.\\
\end{longtable}

\uctitle{UC-ADM-05}{Quản Trị Viên Xem Dashboard/Nhật Ký}
\begin{longtable}{>{\raggedright\arraybackslash}p{0.28\linewidth}>{\raggedright\arraybackslash}p{0.66\linewidth}}
Bằng chứng & \texttt{getDashboard}, \texttt{getLogs}, \texttt{AdminDashboardPage.jsx}.\\
Ghi chú & RB-12: log chưa đầy đủ mọi hành động admin.\\
\end{longtable}

\uctitle{UC-ADM-06}{Quản Trị Viên Quản Lý Người Dùng}
\begin{longtable}{>{\raggedright\arraybackslash}p{0.28\linewidth}>{\raggedright\arraybackslash}p{0.66\linewidth}}
Actor chính & Quản trị viên.\\
Mục tiêu & Xem danh sách người dùng, khóa/mở tài khoản, thay đổi vai trò khi cần.\\
Luồng chính & 1. GET \texttt{/api/admin/users}. 2. PATCH \texttt{/api/admin/users/:id/status}. 3. PATCH \texttt{/api/admin/users/:id/role}.\\
Ngoại lệ & E1: User không tồn tại hoặc role không hợp lệ.\\
Bằng chứng & \texttt{getUsers}, \texttt{updateUserStatus}, \texttt{updateUserRole}, \texttt{AdminDashboardPage.jsx}.\\
\end{longtable}

\uctitle{UC-ADM-07}{Quản Trị Viên Quản Lý Nội Dung}
\begin{longtable}{>{\raggedright\arraybackslash}p{0.28\linewidth}>{\raggedright\arraybackslash}p{0.66\linewidth}}
Actor chính & Quản trị viên.\\
Mục tiêu & Quản lý danh mục, banner, trang nội dung và popup hiển thị ở frontend.\\
Luồng chính & Sử dụng các API CRUD \texttt{/api/admin/content/categories}, \texttt{/banners}, \texttt{/pages}, \texttt{/popups}.\\
Hậu điều kiện & Nội dung public được frontend lấy qua \texttt{content.service.js}.\\
Bằng chứng & \texttt{admin.controller.js}, \texttt{admin.routes.js}, \texttt{AdminDashboardPage.jsx}, \texttt{ContentPage.jsx}.\\
\end{longtable}

\uctitle{UC-ADM-08}{Quản Trị Viên Quản Lý Đối Tác Mở Rộng}
\begin{longtable}{>{\raggedright\arraybackslash}p{0.28\linewidth}>{\raggedright\arraybackslash}p{0.66\linewidth}}
Actor chính & Quản trị viên.\\
Mục tiêu & Quản lý trạng thái đối tác, chi nhánh đối tác và các yêu cầu mở khóa.\\
Luồng chính & 1. PATCH \texttt{/api/admin/partners/:id/status}. 2. GET \texttt{/api/admin/partners/:id/branches}. 3. PATCH \texttt{/api/admin/partners/branches/:id}. 4. GET/PATCH \texttt{/api/admin/partner-appeals}.\\
Ngoại lệ & E1: Partner/branch/appeal không tồn tại hoặc trạng thái cập nhật không hợp lệ.\\
Bằng chứng & \texttt{updatePartnerStatusAny}, \texttt{getPartnerBranches}, \texttt{updatePartnerBranchStatus}, \texttt{updatePartnerAppealStatus}.\\
\end{longtable}

\section{Luồng Ngoại Lệ Chung}
Trùng email, voucher chưa duyệt, hết stock, thanh toán thất bại, scan sai đối tác, review không thuộc khách — xử lý qua \texttt{BusinessException} và \texttt{error.middleware.js}.


\chapter{ERD Và Mô Hình Dữ Liệu}
\section{Tổng Quan}
Schema PostgreSQL định nghĩa tại \texttt{backend/src/config/init.sql}. Dữ liệu mẫu bổ sung trong \texttt{seed-data.sql}.

\textbf{NOTE: Nhóm cần xuất ERD từ database hoặc vẽ lại bằng draw.io rồi chèn hình tại đây.}

\section{Bảng Danh Sách Entity}
\small
\begin{longtable}{>{\raggedright\arraybackslash}p{0.22\linewidth}>{\raggedright\arraybackslash}p{0.50\linewidth}>{\raggedright\arraybackslash}p{0.20\linewidth}}
\toprule
\textbf{Entity/Bảng} & \textbf{Mô tả nghiệp vụ} & \textbf{File schema}\\
\midrule
\endhead
users & Tài khoản đăng nhập, phân quyền & \texttt{init.sql} L22--32\\
partners & Hồ sơ doanh nghiệp đối tác & L35--46\\
partner\_branches & Chi nhánh áp dụng voucher & L49--57\\
vouchers & Sản phẩm voucher & L60--82\\
voucher\_applicable\_branches & N-N voucher--chi nhánh & L85--90\\
orders & Đơn hàng khách & L93--107\\
order\_items & Chi tiết đơn & L119--127\\
issued\_vouchers & Mã phát hành & L131--143\\
reviews & Đánh giá & L147--158\\
complaints & Khiếu nại khách & L161--175\\
partner\_appeals & Khiếu nại đối tác bị khóa & L178--191\\
categories, banners, content\_pages, popups & Quản trị nội dung & L194--232\\
system\_logs & Nhật ký hệ thống & L235--244\\
\bottomrule
\end{longtable}
\normalsize

\section{Bảng Khóa Chính Và Khóa Ngoại}
\small
\begin{longtable}{>{\raggedright\arraybackslash}p{0.20\linewidth}>{\raggedright\arraybackslash}p{0.12\linewidth}>{\raggedright\arraybackslash}p{0.58\linewidth}}
\toprule
\textbf{Bảng} & \textbf{PK} & \textbf{FK chính}\\
\midrule
\endhead
users & id (UUID) & —\\
partners & id & user\_id → users(id)\\
partner\_branches & id & partner\_id → partners(id)\\
vouchers & id & partner\_id → partners(id)\\
orders & id & customer\_id → users(id)\\
order\_items & id & order\_id → orders(id); voucher\_id → vouchers(id)\\
issued\_vouchers & id & order\_item\_id, voucher\_id, customer\_id, partner\_id; used\_at\_branch → partner\_branches\\
reviews & id & voucher\_id, customer\_id, issued\_voucher\_id\\
complaints & id & customer\_id; optional voucher\_id, issued\_voucher\_id, order\_id\\
partner\_appeals & id & partner\_id, user\_id\\
system\_logs & id (BIGSERIAL) & user\_id → users(id) nullable\\
\bottomrule
\end{longtable}
\normalsize

\section{Bảng Quan Hệ Giữa Các Entity}
\begin{longtable}{>{\raggedright\arraybackslash}p{0.18\linewidth}>{\raggedright\arraybackslash}p{0.10\linewidth}>{\raggedright\arraybackslash}p{0.18\linewidth}>{\raggedright\arraybackslash}p{0.46\linewidth}}
\toprule
\textbf{Entity A} & \textbf{QH} & \textbf{Entity B} & \textbf{Ghi chú}\\
\midrule
users & 1--1/1--n & partners & Một user PARTNER có một partner profile.\\
partners & 1--n & partner\_branches & Chi nhánh vật lý.\\
partners & 1--n & vouchers & Đối tác sở hữu voucher.\\
vouchers & n--n & partner\_branches & Qua \texttt{voucher\_applicable\_branches}.\\
users & 1--n & orders & Khách đặt hàng.\\
orders & 1--n & order\_items & Chi tiết mua.\\
order\_items & 1--n & issued\_vouchers & Mỗi quantity có thể nhiều mã.\\
issued\_vouchers & 1--1 & reviews & UNIQUE issued\_voucher\_id.\\
users & 1--n & complaints & Khiếu nại.\\
\bottomrule
\end{longtable}

\section{Mapping Yêu Cầu Dữ Liệu Nghiệp Vụ}
\begin{longtable}{>{\raggedright\arraybackslash}p{0.14\linewidth}>{\raggedright\arraybackslash}p{0.30\linewidth}>{\raggedright\arraybackslash}p{0.48\linewidth}}
\toprule
\textbf{Yêu cầu} & \textbf{Dữ liệu nghiệp vụ} & \textbf{Bảng/trường}\\
\midrule
BR-01/RB-01 & Trạng thái duyệt voucher & \texttt{vouchers.status}\\
BR-04/RB-15 & Tồn kho & \texttt{vouchers.stock}\\
BR-05/RB-05 & Mã sau thanh toán & \texttt{issued\_vouchers} sau \texttt{orders.status=PAID}\\
BR-05/RB-07 & Vòng đời mã & \texttt{issued\_vouchers.status}, \texttt{used\_at}\\
BR-CUS-08 & Đánh giá & \texttt{reviews}\\
BR-ADM-07/RB-12 & Audit & \texttt{system\_logs}\\
\bottomrule
\end{longtable}

\section{Mã Mermaid Gợi Ý (Chuyển Thành Hình)}
\begin{verbatim}
erDiagram
  users ||--o| partners : has
  partners ||--{ partner_branches : owns
  partners ||--{ vouchers : creates
  vouchers }o--o{ partner_branches : applicable
  users ||--{ orders : places
  orders ||--{ order_items : contains
  order_items ||--{ issued_vouchers : generates
  issued_vouchers ||--o| reviews : has
  users ||--{ complaints : files
\end{verbatim}


\chapter{Biểu Đồ Xử Lý Nghiệp Vụ}
\section{Tổng Quan}
Các luồng xử lý được mô tả theo template thống nhất. Sơ đồ hình: \textbf{NOTE: Nhóm cần bổ sung Activity Diagram/BPMN dạng hình cho từng luồng tương ứng.}

\newcommand{\flowblock}[8]{
\subsection{#1}
\begin{longtable}{>{\raggedright\arraybackslash}p{0.26\linewidth}>{\raggedright\arraybackslash}p{0.68\linewidth}}
Mục tiêu & #2\\
Actor & #3\\
Điều kiện bắt đầu & #4\\
Các bước xử lý & #5\\
Điều kiện kết thúc & #6\\
Ngoại lệ & #7\\
Bằng chứng codebase & #8\\
Ghi chú hình & \textbf{NOTE: Nhóm cần bổ sung Activity Diagram/BPMN dạng hình cho luồng này.}\\
\end{longtable}
}

\flowblock{Luồng Đăng Ký/Đăng Nhập}{Xác thực người dùng theo vai trò}{Người dùng}{Chưa có phiên hợp lệ}{1. Nhập credential. 2. API auth validate. 3. Hash/so khớp password. 4. Phát JWT. 5. Lưu AuthContext.}{Token hợp lệ, redirect dashboard}{Trùng email; sai password; 401}{\texttt{auth.controller.js}, \texttt{AuthContext.jsx}}

\flowblock{Luồng Mua Voucher}{Hoàn tất mua voucher hợp lệ}{Khách hàng}{Đã đăng nhập CUSTOMER; giỏ không rỗng}{1. Thêm giỏ. 2. createOrder kiểm RB-01/04/15. 3. Tạo PENDING order.}{Order PENDING tạo thành công}{Voucher chưa duyệt; hết stock}{\texttt{CartPage.jsx}, \texttt{order.controller.js}}

\flowblock{Luồng Thanh Toán}{Chuyển đơn sang PAID}{Khách hàng}{Order PENDING}{1. Chọn COD/VNPay/PayPal. 2. payOrder hoặc callback. 3. Giảm stock. 4. Sinh issued\_vouchers.}{Order PAID, có mã}{Thanh toán thất bại; stock race}{\texttt{payOrder}, \texttt{vnpay.js}, \texttt{paypal.js}}

\flowblock{Luồng Phát Hành Voucher Code}{Sinh mã duy nhất sau thanh toán}{Hệ thống}{Order chuyển PAID}{1. Duyệt order\_items. 2. generateCode VCH-*. 3. INSERT issued\_vouchers với expires\_at.}{Bản ghi UNUSED trong DB}{Đơn CANCELLED không sinh mã}{RB-05: \texttt{order.controller.js}, bảng \texttt{issued\_vouchers}}

\flowblock{Luồng Sử Dụng Voucher}{Xác nhận khách đã dùng dịch vụ}{Đối tác}{Partner APPROVED; mã hợp lệ}{1. Nhập code. 2. checkVoucher/scanVoucher. 3. Kiểm partner\_id, branch, hạn. 4. UPDATE USED.}{status=USED}{Đã USED; sai đối tác; hết hạn}{\texttt{partner.controller.js}, \texttt{PartnerVoucherScan.jsx}}

\flowblock{Luồng Đối Tác Tạo Và Gửi Duyệt}{Tạo voucher chờ admin}{Đối tác}{Partner APPROVED}{1. Form nhập dữ liệu. 2. Validate RB-02/03. 3. Lưu DRAFT hoặc PENDING\_APPROVAL. 4. Gắn chi nhánh.}{Voucher lưu DB}{Partner chưa duyệt bị chặn}{\texttt{PartnerVoucherForm.jsx}, \texttt{voucher.controller.js}}

\flowblock{Luồng Admin Duyệt Voucher}{Kiểm soát voucher được bán}{Admin}{Đăng nhập ADMIN}{1. Lấy pending list. 2. Xem chi tiết. 3. approve/reject. 4. Cập nhật status.}{APPROVED hoặc REJECTED}{—}{\texttt{admin.controller.js}, \texttt{AdminVoucherReview.jsx}}

\flowblock{Luồng Admin Duyệt Đối Tác}{Kích hoạt đối tác}{Admin}{Partner PENDING}{1. Xem hồ sơ. 2. approve/reject. 3. Cập nhật partners.status.}{APPROVED — cho phép tạo voucher}{—}{\texttt{approvePartner}, \texttt{AdminDashboardPage.jsx}}

\flowblock{Luồng Khiếu Nại}{Xử lý phản hồi khách}{Khách hàng, Admin}{Khách đã mua/có vấn đề}{1. POST complaint. 2. Admin xem list. 3. PATCH status + response.}{RESOLVED/REJECTED}{—}{\texttt{user.controller.js}, \texttt{admin.controller.js}, \texttt{complaints}}

\flowblock{Luồng Đánh Giá Voucher}{Ghi nhận đánh giá sau khi mua}{Khách hàng, Đối tác}{Khách đã có issued voucher hợp lệ}{1. Khách chọn mã đã mua. 2. POST review. 3. Đối tác có thể reply.}{Review lưu DB, có thể có phản hồi}{Review không thuộc mã của khách; trùng lỗi}{\texttt{review.controller.js}, \texttt{VoucherDetailPage.jsx}, \texttt{review.routes.js}}

\flowblock{Luồng Báo Cáo/Bảng Điều Khiển}{Tổng hợp thống kê vận hành}{Admin, Đối tác}{Đã đăng nhập}{1. Gọi dashboard/report APIs. 2. Aggregate số liệu. 3. Hiển thị KPI.}{Số liệu dashboard/báo cáo}{—}{\texttt{admin.controller.js}, \texttt{partner.controller.js}, \texttt{AdminDashboardPage.jsx}, \texttt{PartnerReports.jsx}}

\flowblock{Luồng Đối Tác Gửi Yêu Cầu Mở Khóa}{Gửi appeal khi bị tạm khóa}{Đối tác, Admin}{Partner SUSPENDED}{1. Nhập tiêu đề, nội dung. 2. POST appeal. 3. Admin xem và cập nhật.}{Appeal được ghi nhận}{Chưa bị khóa; đã có appeal PENDING}{\texttt{partner.controller.js}, \texttt{admin.controller.js}, \texttt{partner\_appeals}}


\chapter{Từ Điển Dữ Liệu}
\section{Giới Thiệu}
Từ điển dữ liệu mô tả các bảng quan trọng theo \texttt{backend/src/config/init.sql}. Kiểu dữ liệu lấy từ schema; nếu không rõ ghi \textit{Cần kiểm tra thêm}.

\section{Quy Ước}
Snake\_case; PK thường UUID (trừ \texttt{system\_logs.id} BIGSERIAL); enum PostgreSQL cho trạng thái.

\newcommand{\dicttable}[1]{\subsection{#1}}
\newcommand{\dicthead}{\begin{longtable}{>{\raggedright\arraybackslash}p{0.16\linewidth}>{\raggedright\arraybackslash}p{0.12\linewidth}>{\raggedright\arraybackslash}p{0.08\linewidth}>{\raggedright\arraybackslash}p{0.10\linewidth}>{\raggedright\arraybackslash}p{0.22\linewidth}>{\raggedright\arraybackslash}p{0.14\linewidth}>{\raggedright\arraybackslash}p{0.10\linewidth}}
\toprule
\textbf{Trường} & \textbf{Kiểu} & \textbf{BB} & \textbf{Khóa} & \textbf{Mô tả} & \textbf{Quy tắc} & \textbf{GC}\\
\midrule
\endhead}
\newcommand{\dictfoot}{\bottomrule\end{longtable}}

\dicttable{users}
\small\dicthead
id & UUID & Có & PK & Định danh & DEFAULT gen\_random\_uuid() & \\
email & VARCHAR(255) & Không & UQ & Email đăng nhập & Nullable cho phone-only & \\
password & VARCHAR(255) & Có & — & Hash bcrypt & Không lưu plain text & \\
full\_name & VARCHAR(255) & Có & — & Họ tên & & \\
phone & VARCHAR(20) & Không & UQ* & SĐT & Partial unique index & \\
role & user\_role & Có & — & ADMIN/PARTNER/CUSTOMER & DEFAULT CUSTOMER & \\
is\_active & BOOLEAN & Có & — & Khóa tài khoản & DEFAULT TRUE & \\
created\_at & TIMESTAMPTZ & Có & — & Tạo & DEFAULT NOW() & \\
updated\_at & TIMESTAMPTZ & Có & — & Cập nhật & & \\
\dictfoot\normalsize

\dicttable{partners}
\small\dicthead
id & UUID & Có & PK & Đối tác & & \\
user\_id & UUID & Có & FK & Liên kết user & ON DELETE CASCADE & \\
business\_name & VARCHAR(255) & Có & — & Tên DN & & \\
business\_license & VARCHAR(100) & Không & — & Giấy phép & & \\
representative & VARCHAR(255) & Có & — & Người đại diện & & \\
address & TEXT & Không & — & Địa chỉ & & \\
status & partner\_status & Có & — & PENDING/APPROVED/... & DEFAULT PENDING & \\
rejection\_reason & TEXT & Không & — & Lý do từ chối & & \\
\dictfoot\normalsize

\dicttable{partner\_branches}
\small\dicthead
id & UUID & Có & PK & Chi nhánh & & \\
partner\_id & UUID & Có & FK & Thuộc đối tác & CASCADE & \\
name & VARCHAR(255) & Có & — & Tên chi nhánh & & \\
address & TEXT & Có & — & Địa chỉ & & \\
phone & VARCHAR(20) & Không & — & Liên hệ & & \\
is\_active & BOOLEAN & Có & — & Hoạt động & DEFAULT TRUE & \\
\dictfoot\normalsize

\dicttable{vouchers}
\small\dicthead
id & UUID & Có & PK & Voucher & & \\
partner\_id & UUID & Có & FK & Chủ sở hữu & & \\
name & VARCHAR(255) & Có & — & Tên SP & & \\
original\_price & NUMERIC(15,2) & Có & — & Giá gốc & RB-02: $>$ sale\_price & \\
sale\_price & NUMERIC(15,2) & Có & — & Giá bán & chk\_price & \\
stock & INTEGER & Có & — & Tồn & $\geq 0$; RB-15 & \\
sale\_start/end & TIMESTAMPTZ & Không & — & Thời gian bán & RB-03 & \\
valid\_until & TIMESTAMPTZ & Không & — & Hạn dùng sau mua & & \\
status & voucher\_status & Có & — & Vòng đời & RB-01 & \\
\dictfoot\normalsize

\dicttable{orders}
\small\dicthead
id & UUID & Có & PK & Đơn hàng & & \\
customer\_id & UUID & Có & FK & Khách & & \\
total\_amount & NUMERIC(15,2) & Có & — & Tổng tiền & & \\
status & order\_status & Có & — & PENDING/PAID/... & RB-13 & \\
payment\_method & VARCHAR(50) & Không & — & COD/VNPAY/PAYPAL & & \\
payment\_ref & VARCHAR(255) & Không & — & Mã GD & & \\
recipient\_* & VARCHAR & Không & — & Người nhận quà & & \\
paid\_at & TIMESTAMPTZ & Không & — & Thời điểm TT & & \\
\dictfoot\normalsize

\dicttable{order\_items}
\small\dicthead
id & UUID & Có & PK & Dòng đơn & & \\
order\_id & UUID & Có & FK & Thuộc order & CASCADE & \\
voucher\_id & UUID & Có & FK & Voucher mua & & \\
quantity & INTEGER & Có & — & SL & $>$ 0 & \\
unit\_price & NUMERIC(15,2) & Có & — & Giá tại thời điểm mua & & \\
\dictfoot\normalsize

\dicttable{issued\_vouchers}
\small\dicthead
code & VARCHAR(64) & Có & UQ & Mã voucher & RB-06 unique & \\
order\_item\_id & UUID & Có & FK & Nguồn đơn & RB-05 & \\
status & issued\_voucher\_status & Có & — & UNUSED/USED/... & RB-07 & \\
used\_at & TIMESTAMPTZ & Không & — & Thời điểm dùng & & \\
used\_at\_branch & UUID & Không & FK & Chi nhánh dùng & RB-09 & \\
expires\_at & TIMESTAMPTZ & Không & — & Hết hạn mã & RB-08 & \\
\dictfoot\normalsize

\dicttable{reviews}
\small\dicthead
issued\_voucher\_id & UUID & Có & FK,UQ & Một review/mã & RB-10 & \\
rating & SMALLINT & Có & — & 1--5 sao & CHECK 1--5 & \\
comment & TEXT & Không & — & Bình luận & & \\
partner\_reply & TEXT & Không & — & Trả lời đối tác & & \\
\dictfoot\normalsize

\dicttable{complaints}
\small\dicthead
customer\_id & UUID & Có & FK & Người gửi & & \\
subject & VARCHAR(255) & Có & — & Tiêu đề & & \\
message & TEXT & Có & — & Nội dung & & \\
status & complaint\_status & Có & — & PENDING/... & & \\
admin\_response & TEXT & Không & — & Phản hồi admin & & \\
\dictfoot\normalsize

\dicttable{partner\_appeals}
\small\dicthead
partner\_id & UUID & Có & FK & Đối tác khiếu nại & & \\
title & VARCHAR(255) & Có & — & Tiêu đề & & \\
content & TEXT & Có & — & Nội dung & & \\
status & partner\_appeal\_status & Có & — & PENDING/... & 1 pending/partner & \\
\dictfoot\normalsize

\dicttable{system\_logs}
\small\dicthead
id & BIGSERIAL & Có & PK & ID log & & \\
action & VARCHAR(100) & Có & — & Hành động & RB-12 & \\
entity & VARCHAR(100) & Không & — & Loại entity & & \\
details & JSONB & Không & — & Chi tiết & & \\
ip\_address & VARCHAR(45) & Không & — & IP client & & \\
\dictfoot\normalsize

\section{Ghi Chú}
\textbf{NOTE: Nhóm cần đối chiếu lại file schema/migration để bổ sung kiểu dữ liệu chính xác nếu schema thay đổi sau này.}


\chapter{Đặc Tả Thiết Kế Giao Diện}
\section{Tổng Quan}
Giao diện chia theo vai trò; routing tại \texttt{frontend/src/App.jsx}. Nguyên tắc: rõ luồng mua hàng, hiển thị trạng thái, phân quyền qua \texttt{ProtectedRoute}.

\section{Màn Hình Khách Hàng}
\small
\begin{longtable}{>{\raggedright\arraybackslash}p{0.14\linewidth}>{\raggedright\arraybackslash}p{0.12\linewidth}>{\raggedright\arraybackslash}p{0.52\linewidth}>{\raggedright\arraybackslash}p{0.14\linewidth}}
\toprule
\textbf{Màn hình} & \textbf{Vai trò} & \textbf{Thành phần / Hành động / File} & \textbf{Cải thiện}\\
\midrule
\endhead
Trang chủ & CUSTOMER, public & \texttt{HomePage.jsx}: banner, danh mục, flash sale, popup & Thêm ảnh thật\\
Danh sách voucher & CUSTOMER, public & \texttt{VouchersPage.jsx}: lọc, sort, phân trang & Lọc khu vực\\
Chi tiết voucher & CUSTOMER, public & \texttt{VoucherDetailPage.jsx}: giá, chi nhánh, thêm giỏ, review & —\\
Giỏ hàng & CUSTOMER & \texttt{CartPage.jsx}: sửa SL, người nhận, chọn TT & Validate client\\
Thanh toán & CUSTOMER & Tích hợp trong Cart; VNPay \texttt{VnpayReturnPage.jsx} & —\\
Voucher của tôi & CUSTOMER & \texttt{MyVouchersPage.jsx}: đơn, mã, QR mô phỏng & —\\
Đánh giá/khiếu nại & CUSTOMER & Form trên Orders/MyVouchers; \texttt{complaint.service.js} & —\\
Hồ sơ cá nhân & All auth & \texttt{ProfilePage.jsx}: sửa profile, đổi MK & —\\
\bottomrule
\end{longtable}
\normalsize

\section{Màn Hình Đối Tác}
\small
\begin{longtable}{>{\raggedright\arraybackslash}p{0.16\linewidth}>{\raggedright\arraybackslash}p{0.12\linewidth}>{\raggedright\arraybackslash}p{0.52\linewidth}>{\raggedright\arraybackslash}p{0.12\linewidth}}
\toprule
\textbf{Màn hình} & \textbf{Vai trò} & \textbf{Chi tiết} & \textbf{GC}\\
\midrule
Dashboard đối tác & PARTNER & \texttt{PartnerDashboardPage.jsx}: KPI, trạng thái, appeal & —\\
Hồ sơ đối tác & PARTNER & Tab hồ sơ, \texttt{updateProfile} & —\\
Quản lý chi nhánh & PARTNER & CRUD branch trong dashboard & —\\
Tạo/sửa voucher & PARTNER & \texttt{PartnerVoucherForm.jsx} & —\\
Danh sách voucher & PARTNER & \texttt{PartnerVouchers.jsx} & —\\
Quét/nhập mã & PARTNER & \texttt{PartnerVoucherScan.jsx} & QR camera: CT thêm\\
Báo cáo & PARTNER & \texttt{PartnerReports.jsx} & —\\
Khiếu nại đối tác & PARTNER SUSPENDED & Appeal form, \texttt{partner\_appeals} & —\\
\bottomrule
\end{longtable}
\normalsize

\section{Màn Hình Quản Trị Viên}
\small
\begin{longtable}{>{\raggedright\arraybackslash}p{0.16\linewidth}>{\raggedright\arraybackslash}p{0.12\linewidth}>{\raggedright\arraybackslash}p{0.52\linewidth}>{\raggedright\arraybackslash}p{0.12\linewidth}}
\toprule
\textbf{Màn hình} & \textbf{Vai trò} & \textbf{Chi tiết} & \textbf{GC}\\
\midrule
Dashboard QT & ADMIN & \texttt{AdminDashboardPage.jsx}: tab tổng hợp & Tách module\\
QL người dùng & ADMIN & Tab users, PATCH status/role & —\\
QL đối tác & ADMIN & Duyệt/từ chối/khóa partner & —\\
Duyệt voucher & ADMIN & \texttt{AdminVoucherReview.jsx} + tab voucher & —\\
QL đơn hàng & ADMIN & Tab orders, cập nhật status & —\\
QL nội dung & ADMIN & categories/banners/pages/popups & —\\
Xử lý khiếu nại & ADMIN & Tab complaints + partner appeals & —\\
Nhật ký HT & ADMIN & Tab logs, \texttt{getLogs} & —\\
\bottomrule
\end{longtable}
\normalsize

\section{Điều Kiện Hiển Thị/Chặn}
\begin{itemize}[leftmargin=1.2cm]
  \item Route CUSTOMER: \texttt{ProtectedRoute roles=["CUSTOMER"]}.
  \item Partner chưa APPROVED: \texttt{usePartnerStatus.js} chặn tạo voucher.
  \item Admin/partner không vào \texttt{/cart}: redirect \texttt{/unauthorized}.
\end{itemize}


\chapter{Kết Luận}
\section{Tóm Tắt Mức Độ Hoàn Thành}
Hệ thống EC-Voucher đã triển khai phần lớn yêu cầu trọng tâm của BRD: ba vai trò người dùng, quản lý voucher có duyệt, mua hàng trực tuyến, thanh toán mô phỏng (COD/VNPay/PayPal), phát hành mã sau thanh toán, xác nhận sử dụng tại đối tác, quản trị nội dung và dashboard. Ma trận đối chiếu (Chương Đối Chiếu Yêu Cầu) cho thấy đa số BR/RB ở mức \textit{Hoàn thành} hoặc \textit{Một phần} — không khẳng định đáp ứng 100\% toàn bộ yêu cầu.

\section{Luồng Chính Đã Có}
Đăng ký/đăng nhập/phân quyền; đối tác tạo và gửi duyệt voucher; admin duyệt đối tác/voucher; khách tìm kiếm, giỏ hàng, đặt hàng, thanh toán; sinh \texttt{issued\_vouchers}; đối tác kiểm tra/xác nhận mã; đánh giá và khiếu nại; báo cáo dashboard.

\section{Hạn Chế Và Việc Cần Hoàn Thiện}
\begin{itemize}[leftmargin=1.2cm]
  \item \textbf{Audit log (RB-12/BR-ADM-07)}: Bảng \texttt{system\_logs} có nhưng chưa ghi đủ mọi thao tác quản trị.
  \item \textbf{Vòng đời voucher/code}: Chưa có scheduled job tự động EXPIRED — \textbf{Chưa tìm thấy trong codebase}.
  \item \textbf{Hoàn tiền (RB-14)}: Chỉ cập nhật thủ công trạng thái REFUNDED.
  \item \textbf{Kiểm thử}: Test plan có 25+ case nhưng phần lớn \textit{Chưa chạy}; cần cập nhật kết quả thực tế.
  \item \textbf{Tài liệu hình ảnh}: ERD, BPMN/Activity Diagram, screenshot UI — nhóm cần bổ sung thủ công.
  \item \textbf{Module admin}: \texttt{AdminDashboardPage.jsx} tập trung nhiều chức năng, nên tách nhỏ.
\end{itemize}

\section{Ưu Tiên Trước Khi Nộp}
\begin{enumerate}[leftmargin=1.5cm]
  \item Sửa và xác nhận ma trận đối chiếu BR/RB (đã tách đúng trong báo cáo này).
  \item Chạy test case bắt buộc (TC-A01--A04, TC-C03--C08, TC-P04--P06, TC-AD01--AD02, TC-R01--R02).
  \item Chèn logo, ảnh UI, ERD, slide, link repo/video demo.
  \item Rà soát các mục \textit{Cần kiểm tra thêm} trong bảng đối chiếu.
\end{enumerate}

\section{Kết Luận Chung}
Đồ án đạt được mục tiêu xây dựng nền tảng bán voucher có kiểm soát nghiệp vụ và bằng chứng triển khai rõ trong codebase. Hệ thống phù hợp demo và học tập; để đưa vào vận hành thực tế cần hoàn thiện kiểm toán, kiểm thử, tự động hóa vòng đời voucher và tách module quản trị.



\end{document}


